\documentclass[UTF8,a4paper,11pt]{ctexart}
\usepackage[margin=2.15cm]{geometry}
\usepackage{booktabs}
\usepackage{longtable}
\usepackage{array}
\usepackage{xcolor}
\usepackage{hyperref}
\hypersetup{colorlinks=true,linkcolor=blue!45!black,urlcolor=blue!45!black}
\setlength{\emergencystretch}{3em}
\setlength{\parskip}{0.35em}
\setlength{\parindent}{2em}

\title{\heiti WealthCraftAgent 股票分析 Agent 反思报告}
\author{2352048 夏浩博 \quad 2351269 黄一和}
\date{\today}

\begin{document}
\maketitle

\section{项目目标与领域选择}
本项目以 WealthCraftAgent 股票分析 Agent 作为 .NET AI Agent 期末项目主题。股票分析天然需要多源信息融合：公司主营业务帮助理解企业基本面，K 线帮助观察价格结构，财务报表反映经营质量，新闻动态反映事件冲击和行业变化。因此该领域很适合作为 Agent 项目：模型不能只聊天，还必须能规划任务、调用工具、读取数据、维护上下文并生成结构化分析。

项目最终实现了一个基于 .NET 10 和 WPF 的桌面端应用。用户输入股票代码或股票名称后，系统自动创建分析会话，并从主营业务、K 线、新闻、财务四个角度形成报告。在同一股票会话内，系统支持“只看消极新闻”“用缠论分析最近走势”“看五年财务”等追问，并根据意图只调度相关子 Agent。

\section{Agent 内部工作原理}
课程要求的 Agent 核心是 LLM、Agent Loop、Memory 与 Tools。项目中这四个组件分别对应：

\begin{itemize}
  \item LLM：通过 Microsoft Semantic Kernel 接入 OpenAI 兼容 ChatCompletion 接口。
  \item Agent Loop：\texttt{InvestAgentLoop} 实现标准 Thought--Action--Observation 循环；会话编排器实现面向股票分析的多 Agent 工作流。
  \item Memory：对话记忆保存系统提示、用户消息、助手回复和工具返回；会话状态保存当前股票的结构化数据。
  \item Tools：通过 Kernel 插件注册股票价格、财务、新闻、技术指标、本地知识库等函数。
\end{itemize}

项目采用“通用 ReAct + 业务工作流”的混合方案。通用 ReAct 对应课程中的 Agent Loop 原理；业务工作流用于提升演示与运行稳定性。股票分析具有固定模块，因此系统由 Agent A 先规划任务，再并行调度 B/C/D，减少纯自由工具调用带来的不确定性。

\section{关键设计决策}
\subsection{为什么使用 Semantic Kernel}
直接 HTTP 调用 LLM 也能完成对话，但工具调用、消息历史、流式输出、图像输入和函数注册都需要额外维护。Semantic Kernel 能把 C\# 方法映射成模型可调用的 Kernel Function，并提供 \texttt{FunctionChoiceBehavior.Auto()}，更贴近课程中“Tool Calling”的要求，也体现了 .NET 技术栈在 Agent 工程中的应用深度。

\subsection{为什么拆分为 Agent A/B/C/D}
股票分析中不同维度的数据来源和提示词策略差异很大。K 线需要价格序列、技术指标、缠论与历史走势；新闻需要事件筛选和情绪分类；财务需要多年指标趋势。拆分后每个子 Agent 的输入、工具、提示词和输出都更清晰，Agent A 只负责协调、派单和总结，符合单一职责原则。

\subsection{为什么加入 RAG}
缠论和历史相似走势都不是通用模型稳定掌握的项目私有知识。缠论 RAG 从本地 Markdown 与图片索引检索概念、模板和图例；历史趋势 RAG 从本地案例库匹配相似结构，提供样本数量和后续表现。RAG 机制使模型回答具有明确资料依据，同时在输出中说明“不确定性”和“历史不代表预测”。

\section{多 Agent 协作、RAG 与记忆机制}
首轮分析时，Agent A 先补充主营业务，再创建三个任务：Agent B 分析 K 线，Agent C 分析新闻，Agent D 分析财务。三个子任务通过 \texttt{Task.WhenAll} 并发执行，结果写回 \texttt{AnalysisSessionState}，最后 Agent A 读取 B/C/D 原文并生成最终风险提示。

追问阶段，Agent A 会把用户自然语言转成结构化派单 JSON。例如“只看利空新闻”会派给 Agent C，并设置 \texttt{NewsSentimentFilter=negative}；“用缠论看最近走势”会派给 Agent B，并设置 \texttt{UseChanTheory=true}。如果用户在当前会话内要求切换到另一只股票，系统会拒绝并提示新建会话，避免上下文混淆。

记忆机制分两层。第一层是对话记忆，用于给 LLM 提供最近上下文；第二层是结构化状态，用于保存 K 线缓存、新闻、财务序列和各 Agent 结果。结构化状态比纯文本历史更可靠，UI 展示、历史记录和追问处理均可复用这些字段。

\section{错误处理与流式输出}
项目在三个位置处理错误。第一，LLM 调用失败时记录日志并返回用户可读错误。第二，子 Agent 执行失败时不会让整个分析崩溃，而是记录失败步骤并继续后续模块。第三，数据源不足或知识库未命中时，提示词要求模型明确说明不确定性，避免编造结论。

流式输出由 \texttt{AgentPromptRunner} 控制。它按时间和字符数节流推送 partial 内容，避免 UI 更新过于频繁；同时过滤 \texttt{<think>}、\texttt{analysis} 等推理草稿标签，保证用户看到的是最终可读分析，而不是模型内部推理过程。

\section{测试与不足}
测试主要覆盖 DI 容器、插件注册、Agent Loop 基础结构、组合数据服务和历史走势服务。测试能证明核心服务可解析、插件函数可发现、关键服务的输入输出结构正常。

不足主要有三点：第一，股票数据来源依赖外部服务，网络或接口变化会影响结果稳定性；第二，RAG 检索目前以关键词和特征匹配为主，没有接入真正的向量数据库；第三，股票分析本身具有高度不确定性，系统必须持续强调“仅供学习研究，不构成投资建议”。

\section{AI 使用透明度}
本项目开发中使用了 AI 进行大量辅助。AI 主要参与代码生成、局部重构、文档起草、测试建议和错误排查建议。小组成员负责确定需求范围、检查生成代码是否符合项目目标、运行和调试项目、整合不同模块、核对文档内容，并对核心代码进行逐行理解。AI 在本项目中作为开发效率工具，不替代成员对代码和架构的最终理解。

\section{Agent 核心循环代码逐行解读}
本节解读 \texttt{src/InvestAgent.Core/Agent/InvestAgentLoop.cs} 第 1--177 行。该文件是课程要求中最典型的 ReAct Agent Loop 示例。

\begin{longtable}{p{2.4cm}p{12.6cm}}
\toprule
行号 & 作用解读\\
\midrule
\endhead
1--3 & 引入项目内部命名空间：配置项 \texttt{AgentOptions}、记忆接口 \texttt{IConversationMemory}、步骤模型 \texttt{AgentStep} 等。没有这些类型，循环无法读取最大步数、写入上下文或返回步骤。\\
4 & 引入 Microsoft 日志接口，用于记录 Agent 开始、完成、失败、达到最大步数等运行状态。\\
5--7 & 引入 Semantic Kernel、聊天补全接口和 OpenAI 执行设置。这些类型支撑 LLM 调用和自动工具调用。\\
9 & 声明当前类所在命名空间 \texttt{InvestAgent.Core.Agent}，表明这是 Core 层 Agent 逻辑的一部分。\\
11--16 & XML 注释说明该类是投资分析 Agent 的主执行循环，实现 ReAct 模式，并通过事件向 UI 推送每一步状态。该位置是 Thought--Action--Observation 代码结构的入口说明。\\
17 & 定义公开类 \texttt{InvestAgentLoop}。它不是 UI 控件，也不是数据服务，而是连接 LLM、工具、记忆和步骤流的控制器。\\
18 & 类体开始。\\
19 & 字段 \texttt{\_kernel} 保存 Semantic Kernel 实例。工具函数都注册在 Kernel 中，后续工具调用会通过它执行。\\
20 & 字段 \texttt{\_chatService} 保存聊天补全服务，负责真正向 LLM 发送消息并获取回复。\\
21 & 字段 \texttt{\_memory} 保存对话记忆，循环每轮都从这里读取历史，并把用户、助手、工具消息写回去。\\
22 & 字段 \texttt{\_options} 保存配置，如最大步数。它让循环不会无限执行。\\
23 & 字段 \texttt{\_logger} 保存日志对象，用于记录运行信息和异常。\\
25--26 & 定义 \texttt{OnStep} 事件。每产生一个 Thought、Action、Observation 或 Response，就触发该事件，UI 能够订阅并实时显示 Agent 过程。\\
28--30 & 构造函数签名。通过依赖注入传入 Kernel、LLM 服务、记忆、配置和日志，避免在类内部手动创建依赖。\\
31--35 & 把构造函数参数赋值给私有字段。这样 \texttt{RunAsync} 执行时能够使用这些依赖。\\
36 & 构造函数结束。\\
38--44 & XML 注释说明 \texttt{RunAsync} 是主循环方法，使用异步迭代器，每一步通过 \texttt{yield return} 推送给调用方。\\
45 & 声明 \texttt{RunAsync(string userMessage)}，返回 \texttt{IAsyncEnumerable<AgentStep>}。这意味着调用方支持边执行边接收步骤，而不是等全部结束。\\
46 & 方法体开始。\\
47 & 记录日志，截取用户输入前 100 个字符，避免日志过长。这里用 \texttt{Math.Min} 防止输入不足 100 字符时越界。\\
48 & 把用户消息写入记忆。后续 LLM 会看到这条消息，并基于它决定是否调用工具。\\
50--56 & 创建 OpenAI 执行设置。核心是 \texttt{FunctionChoiceBehavior.Auto()}，表示允许模型自动选择 Kernel 中注册的函数；温度 0.7 控制生成多样性；最大 token 数限制回复长度。\\
58--60 & 开始主循环，步数从 1 到 \texttt{\_options.MaxSteps}。这是防止 Agent 无限工具调用的重要保护。\\
61--69 & 构造 Thought 步骤。第一步提示“分析用户问题”，后续步骤提示“根据已有数据判断是否需要更多信息”。该步骤主要用于可观测性，不直接调用 LLM。\\
70 & 触发 \texttt{OnStep} 事件，把 Thought 推给 UI。\\
71 & 使用 \texttt{yield return} 把 Thought 返回给异步枚举调用方。\\
73--75 & 准备调用 LLM：声明 \texttt{response} 保存模型回复，\texttt{errorMessage} 保存异常后的用户可读错误。\\
76--80 & 在 \texttt{try} 中调用 \texttt{\_chatService.GetChatMessageContentAsync}。输入是记忆中的聊天历史、执行设置和 Kernel。Kernel 被传入后，模型才能看到可调用工具。\\
81--85 & 捕获 LLM 调用异常，写日志，并把异常消息包装成用户可读文本，避免程序直接崩溃。\\
87--99 & 如果 \texttt{errorMessage} 不为空，说明 LLM 调用失败。此时构造 Response 类型的错误步骤，推送给 UI 和调用方，然后 \texttt{yield break} 终止循环。\\
101--102 & 从模型回复的 \texttt{Items} 中筛选 \texttt{FunctionCallContent}。这些对象表示模型请求执行的工具调用。\\
104--118 & 如果没有工具调用，说明模型认为任务已经完成。代码构造最终 Response，把助手回复写入记忆，推送步骤，记录完成日志，并结束循环。\\
120--121 & 如果存在工具调用，进入 Action 和 Observation 阶段。先向记忆中加入一条空助手消息，表示本轮助手已经产生了工具调用意图。\\
123 & 遍历模型请求的每个工具调用。一轮 LLM 回复可能包含多个函数调用。\\
125--133 & 为当前函数调用构造 Action 步骤，记录步骤编号、步骤类型、函数名、参数和展示文本“调用工具”。\\
134--135 & 把 Action 推送给 UI 和异步枚举调用方，让用户知道 Agent 正在调用哪个工具。\\
137--143 & 执行工具函数。关键语句是 \texttt{functionCall.InvokeAsync(\_kernel)}，它会在 Kernel 中找到对应函数并执行。返回值转成字符串；如果为空，就写“工具执行完成”。\\
144--148 & 捕获工具执行异常。工具失败不会让整个循环崩溃，而是记录日志，并把失败原因作为 Observation 返回给模型。\\
150--159 & 构造 Observation 步骤。这里记录函数名和函数返回结果，并把超过 2000 字符的结果截断，避免长 JSON 撑爆上下文或 UI。\\
160--161 & 推送 Observation 给 UI 和调用方。\\
163 & 把完整工具结果写入记忆，角色是 Tool，并带上函数名。下一轮 LLM 会基于这个 Observation 判断是否继续调用工具或输出最终答案。\\
164--165 & 结束当前工具遍历和当前循环步。如果还有步数，回到下一轮 Thought。\\
167--175 & 如果循环跑完仍未结束，说明达到最大步数。代码记录警告，构造最终 Response，提示用户提出更具体的问题，然后推送该步骤。\\
176--177 & 方法和类结束。整个文件实现了从用户输入、LLM 思考、工具行动、结果观察到最终回答的完整闭环。\\
\bottomrule
\end{longtable}

\section{Agent Loop 总结}
Agent Loop 的执行路径为：用户问题先进入 Memory；循环每一步先产生 Thought 以便观察当前状态；随后调用 LLM，模型根据上下文自动选择 Kernel 中注册的工具；当模型请求工具时，系统执行 Action，并将返回结果作为 Observation 写回 Memory；下一轮模型基于 Observation 继续判断；若不再请求工具，则输出最终 Response；若超过最大步数或出现异常，则进入保护性返回路径。

\end{document}
